\DocumentMetadata{
  pdfversion=2.0,
  pdfstandard=ua-2,
  lang=en-US,
  tagging=on,
  tagging-setup={math/setup=mathml-SE,tabsorder=structure}
}
\documentclass[11pt,letterpaper]{article}
\usepackage[margin=0.68in,includefoot]{geometry}
\usepackage{newcomputermodern}
\usepackage{amsmath,amscd,mathtools}
\usepackage{tikz,adjustbox}
\providecommand{\square}{\mdlgwhtsquare}
\usepackage[hidelinks]{hyperref}
\hypersetup{
  pdftitle={Algebra PhD Exam, August 2025},
  pdfauthor={University of Florida Department of Mathematics},
  pdfsubject={Graduate mathematics examination},
  pdfdisplaydoctitle=true,
  bookmarksopen=true
}
\setcounter{secnumdepth}{0}
\setlength{\parindent}{0pt}
\setlength{\parskip}{0.28em}
\setlength{\emergencystretch}{3em}
\setlength{\leftmargini}{2.15em}
\setlength{\leftmarginii}{2.65em}
\setlength{\leftmarginiii}{3em}
\renewcommand{\labelenumi}{\textbf{\theenumi.}}
\pagestyle{empty}
\raggedbottom
\allowdisplaybreaks
\newenvironment{examproblems}[1]{%
  \begin{enumerate}%
  \setcounter{enumi}{#1}%
  \setlength{\topsep}{0.35em}%
  \setlength{\partopsep}{0pt}%
  \setlength{\itemsep}{0.48em}%
  \setlength{\parsep}{0.18em}%
}{\end{enumerate}}
\newenvironment{examsubparts}{%
  \begin{enumerate}%
  \setlength{\topsep}{0.18em}%
  \setlength{\partopsep}{0pt}%
  \setlength{\itemsep}{0.16em}%
  \setlength{\parsep}{0.1em}%
}{\end{enumerate}}
\newenvironment{examinstructions}{%
  \par\begingroup\itshape
}{%
  \par\endgroup\vspace{0.18em}
}
\makeatletter
\renewcommand\section{\@startsection{section}{1}{0pt}%
  {-1.25ex plus -.25ex minus -.1ex}%
  {0.55ex plus .1ex}%
  {\normalfont\large\bfseries}}
\renewcommand{\@maketitle}{%
  \newpage\null\vskip -1em%
  {\footnotesize\bfseries
    \href{https://www.ufl.edu/}{University of Florida}, \href{https://math.ufl.edu/}{Department of Mathematics}\par}%
  \vskip -0.5em%
  \tagstructbegin{tag=Title}%
    {\LARGE\bfseries\href{https://gma.math.ufl.edu/exams/algebra-phd/}{Algebra PhD Exam}, August 2025\par}%
  \tagstructend%
  \vskip 0.25em%
  \tagmcbegin{artifact}%
    \hrule height 1pt%
  \tagmcend%
  \vskip 1em%
}
\makeatother
\title{Algebra PhD Exam, August 2025}
\author{}
\date{}
\begin{document}
\maketitle
\thispagestyle{empty}
\begin{examinstructions}
Answer seven problems. (If you turn in more, the first seven will be graded.) Within reason, you may quote theorems as long as you state them clearly.
\end{examinstructions}
\begin{examinstructions}
Note. Below ring means associative ring with identity, and module means unital module unless otherwise specified.
\end{examinstructions}
\begin{examproblems}{0}
\item
Let~\(K\) be a field and let \(f(x) \in K[x]\) be a non-zero polynomial.
\begin{examsubparts}
\item[\textnormal{(a)}]
Define what is a splitting field for \(f(x)\) over~\(K\).
\item[\textnormal{(b)}]
If \(F_1\) and \(F_2\) are splitting fields for \(f(x)\) over~\(K\) prove, from your definition, that \(F_1\) and \(F_2\) are isomorphic fields.
\end{examsubparts}
\item
Suppose that~\(p\) is an odd prime and the field~\(K\) is the extension of~\(\mathbb{Q}\) generated by a primitive~\(p\)-th root of~\(1\) in~\(\mathbb{C}\).
\begin{examsubparts}
\item[\textnormal{(a)}]
Show that~\(K\) contains a unique quadratic extension~\(F\) of~\(\mathbb{Q}\).
\item[\textnormal{(b)}]
What is \(\operatorname{Gal}(K/F)\)? Justify your answer in detail.
\end{examsubparts}
\item
Let~\(A\) and~\(B\) be finite abelian groups and suppose that \(|A|\) and \(|B|\) are relatively prime. Let \(T=A\otimes_{\mathbb{Z}}B\). Show that \(|T|=1\).
\item
Let~\(\mathcal{C}\) be a concrete category.
\begin{examsubparts}
\item[\textnormal{(a)}]
Define what is meant by a free object in~\(\mathcal{C}\).
\item[\textnormal{(b)}]
Let~\(\mathcal{G}\) be the category of all finite groups. Prove that the free objects in~\(\mathcal{G}\) are exactly the trivial groups.
\end{examsubparts}
\item
Let~\(R\) be a ring with identity. Prove that the following two conditions are equivalent for a left~\(R\)-module~\(P\). (You may not assume any properties of projective modules).
\begin{examsubparts}
\item[\textnormal{1.}]
Every short exact sequence of~\(R\)-modules 
\[
0\to A\to B\to P\to 0
\]
 splits;
\item[\textnormal{2.}]
There is a free~\(R\)-module~\(F\) and an~\(R\)-module~\(K\) such that \(F\simeq K\oplus P\).
\end{examsubparts}
\item
Let~\(K\) be a field and \(K(x)\) the field of rational functions over~\(K\). Let \(\frac{f}{g}\in K(x)\), with \(\frac{f}{g}\notin K\), where~\(f\) and~\(g\) are relatively prime elements of \(K[x]\).
\begin{examsubparts}
\item[\textnormal{(a)}]
Show that \(\frac{f}{g}\) is transcendental over~\(K\).
\item[\textnormal{(b)}]
Prove that \(K(x)\) is a finite extension of \(K\!\left(\frac{f}{g}\right)\).
\end{examsubparts}
\item
State and prove the Lying-Over Theorem on prime ideals and ring extensions.
\item
Prove that an invertible ideal in an integral domain that is a local ring is principal.
\item
Show that for any extension \(A\subseteq B\) of commutative rings with~\(1\), any~\(A\)-module~\(M\) and any~\(B\)-module~\(N\), there is an isomorphism \(\operatorname{Hom}_B(M\otimes_A B,N)\simeq \operatorname{Hom}_A(M,N_A)\) (here \(N_A\) denotes~\(N\) considered as an~\(A\)-module by means of the ring inclusion homomorphism \(A\to B\)).
\item
Let~\(R\) be a commutative ring with~\(1\) with \(1\neq 0\). Prove that the set consisting of~\(0\) and all zero divisors in~\(R\) contains at least one prime ideal of~\(R\).
\item
Classify (up to ring isomorphism) all (not necessarily commutative) semisimple rings of order \(240\).
\end{examproblems}
\end{document}
